\documentclass[10pt]{article}
\providecommand{\FIGDIR}{fig}
\newcommand{\DIGESTDATE}{8 August 2026}
% ===================== PM Research Watch - shared preamble =====================
\usepackage[T1]{fontenc}
\usepackage[utf8]{inputenc}
\usepackage{amsmath,amssymb}
\usepackage{tgpagella}
\usepackage{tgheros}
\renewcommand{\sfdefault}{qhv}
\renewcommand{\ttdefault}{lmtt}
\usepackage{microtype}
\usepackage[english]{babel}

\usepackage[a4paper,top=17mm,bottom=17mm,left=17mm,right=17mm,
            headsep=6mm,footskip=10mm]{geometry}
\usepackage{graphicx}
\usepackage{xcolor}
\usepackage{colortbl}
\usepackage{tikz}
\usetikzlibrary{calc}
\usepackage[most]{tcolorbox}
\usepackage{titlesec}
\usepackage{fancyhdr}
\usepackage{booktabs}
\usepackage{array}
\usepackage{longtable}
\usepackage{enumitem}
\usepackage{multicol}
\usepackage{ragged2e}
\usepackage[hidelinks]{hyperref}

% ------------------------------------------------------------------ palette
\definecolor{Ink}    {HTML}{14262E}
\definecolor{Deep}   {HTML}{0F4C5C}
\definecolor{Teal}   {HTML}{1F7A8C}
\definecolor{Sky}    {HTML}{4FA3C4}
\definecolor{Sage}   {HTML}{6FA287}
\definecolor{Amber}  {HTML}{E3A008}
\definecolor{Clay}   {HTML}{B5643C}
\definecolor{Coral}  {HTML}{D1495B}
\definecolor{Violet} {HTML}{6A5B9E}
\definecolor{Slate}  {HTML}{4A6572}
\definecolor{Paper}  {HTML}{FBF9F5}
\definecolor{Mist}   {HTML}{EEEAE2}
\definecolor{Rule}   {HTML}{D9D3C9}

\pagecolor{Paper}
\color{Ink}

\hypersetup{colorlinks=true, linkcolor=Deep, urlcolor=Teal, citecolor=Teal,
            pdftitle={PM Research Watch}, pdfauthor={Automated daily digest}}

% ------------------------------------------------------------------ headings
\titleformat{\section}
  {\sffamily\bfseries\large\color{Deep}}{}{0pt}
  {}[{\vspace{1.2mm}\color{Amber}\titlerule[1.1pt]}]
\titlespacing*{\section}{0pt}{5.5mm}{2.8mm}

\titleformat{\subsection}
  {\sffamily\bfseries\normalsize\color{Teal}}{}{0pt}{}
\titlespacing*{\subsection}{0pt}{3.6mm}{1.4mm}

\setlist[itemize]{leftmargin=4.2mm, itemsep=1.1mm, topsep=1.3mm, parsep=0pt,
                  label=\textcolor{Teal}{\footnotesize$\blacktriangleright$}}

% ------------------------------------------------------------------ header/footer
\pagestyle{fancy}
\fancyhf{}
\renewcommand{\headrulewidth}{0pt}
\renewcommand{\footrulewidth}{0pt}
\fancyhead[L]{\sffamily\scriptsize\color{Slate}PM RESEARCH WATCH\;\textcolor{Rule}{\textbar}\;\DIGESTDATE}
\fancyhead[R]{\sffamily\scriptsize\color{Slate}Particulate matter monitoring \& health effects}
\fancyfoot[C]{\sffamily\scriptsize\color{Slate}\thepage}

% ------------------------------------------------------------------ components
% metric tile
\newcommand{\metric}[3]{%
  \begin{tcolorbox}[enhanced, colback=white, colframe=Rule, boxrule=0.5pt,
    arc=1.6mm, left=2.2mm, right=2.2mm, top=1.8mm, bottom=1.8mm,
    borderline west={1.7mm}{0pt}{#3}, nobeforeafter, valign=center,
    height=17mm, width=\linewidth]
    \RaggedRight
    {\sffamily\bfseries\fontsize{17}{17}\selectfont\color{#3}#1}\\[0.5mm]
    {\sffamily\fontsize{6.5}{7.4}\selectfont\color{Slate}#2}
  \end{tcolorbox}}

% pull-quote / key insight box
\newtcolorbox{keybox}[1][]{enhanced, colback=Mist, colframe=Mist,
  arc=1.6mm, left=3.2mm, right=3.2mm, top=2.4mm, bottom=2.4mm,
  borderline west={1.7mm}{0pt}{Amber}, #1}

% subtopic band
\newcommand{\band}[2]{%
  \vspace{3.2mm}\noindent
  \begin{tcolorbox}[enhanced, colback=#2, colframe=#2, arc=1.2mm,
    left=2.6mm, right=2.6mm, top=1.1mm, bottom=1.1mm, boxrule=0pt,
    width=\linewidth, nobeforeafter]
    {\sffamily\bfseries\small\color{white}#1}
  \end{tcolorbox}\vspace{1.6mm}}

% one paper entry: {short cite}{title}{journal}{doi}{body}
\newcommand{\paperentry}[5]{%
  \par\vspace{1.9mm}\noindent
  {\sffamily\bfseries\footnotesize\color{Ink}#1}\;%
  {\sffamily\fontsize{7}{8}\selectfont\color{Slate}#3}\\[0.4mm]
  {\footnotesize\itshape\color{Deep}#2}\\[0.6mm]
  {\footnotesize #4}\\[0.3mm]
  {\fontsize{6.6}{7.6}\selectfont\color{Slate}\texttt{doi:}\,\href{https://doi.org/#5}{\texttt{\detokenize{#5}}}}
  \par}

% caption style
\newcommand{\figcap}[1]{\par\vspace{1.0mm}
  {\fontsize{7.2}{8.6}\selectfont\color{Slate}\textsf{\textbf{Fig.}}\; #1\par}\vspace{1.5mm}}

\providecommand{\permil}{\textperthousand}

\setlength{\parindent}{0pt}
\setlength{\parskip}{0pt}
\renewcommand{\arraystretch}{1.18}


\begin{document}

% ==================================================================== MASTHEAD
\thispagestyle{empty}
\begin{tikzpicture}[remember picture, overlay]
  \fill[Deep] (current page.north west) rectangle
      ($(current page.north east)+(0,-67mm)$);
  \fill[Amber] ($(current page.north west)+(0,-67mm)$) rectangle
      ($(current page.north east)+(0,-68.4mm)$);
  \foreach \x/\y/\r in {22/12/0.9, 41/26/1.5, 63/9/0.7, 88/31/1.1, 112/17/1.9,
                        138/29/0.8, 159/13/1.3, 178/54/1.0, 196/21/1.6,
                        30/44/0.6, 74/57/1.2, 125/48/0.9, 168/6/0.8, 191/59/0.7,
                        52/18/1.0, 100/60/0.6, 148/20/1.1, 66/49/1.4, 205/45/0.8}
    \fill[Sky, opacity=0.30] ($(current page.north west)+(\x mm,-\y mm)$) circle (\r mm);
\end{tikzpicture}

\vspace*{4mm}
{\sffamily\fontsize{7.6}{9}\selectfont\color{Sky}%
 AUTOMATED DAILY LITERATURE DIGEST\;\textbullet\;PhD RESEARCH WATCH\par}
\vspace{2.2mm}
{\sffamily\bfseries\fontsize{27}{29}\selectfont\color{white}%
 Particulate Matter\par}
{\sffamily\fontsize{27}{29}\selectfont\color{Sky}%
 Monitoring \& Health Effects\par}
\vspace{5.0mm}
{\sffamily\fontsize{8.4}{10}\selectfont\color{white}%
 Issue 2026-08-08\;\;\textcolor{Amber}{\textbar}\;\;Entry window 8 August 2026\;\;%
 \textcolor{Amber}{\textbar}\;\;14 records in scope\par}
\vspace{18mm}

% -------------------------------------------------------------- metric strip
\noindent
\begin{minipage}[t]{0.190\linewidth}\metric{14}{RECORDS IN SCOPE\\(13 REJECTED)}{Deep}\end{minipage}\hfill
\begin{minipage}[t]{0.190\linewidth}\metric{7}{SENSING / EXPOSURE\\RECORDS (50\%)}{Teal}\end{minipage}\hfill
\begin{minipage}[t]{0.190\linewidth}\metric{7}{SUBTOPIC CLUSTERS\\REPRESENTED}{Sky}\end{minipage}\hfill
\begin{minipage}[t]{0.190\linewidth}\metric{2}{TIER-A\\RECORDS}{Amber}\end{minipage}\hfill
\begin{minipage}[t]{0.190\linewidth}\metric{5}{EFFECT ESTIMATES\\EXTRACTED}{Clay}\end{minipage}

\vspace{6mm}

% -------------------------------------------------------------- signal of day
\section*{\sffamily\bfseries\color{Deep}Signal of the day}
\vspace{-1.5mm}{\color{Amber}\rule{\linewidth}{1.1pt}}\vspace{2.5mm}

\begin{keybox}
{\sffamily\bfseries\small\color{Deep}Two records on opposite ends of the evidence
hierarchy make the same argument: the particle-mass concentration a monitor reports is
conditional on a co-exposure the monitor does not measure.}\\[1.4mm]
\footnotesize
Hur et al.\ (\textit{Environ Pollut}) pooled \textbf{262{,}288} deaths across eight
Kyushu cities over 26 pollen seasons and fitted an explicit SPM\,$\times$\,pollen
interaction in a two-stage time-series. The marginal association of suspended particulate
matter with respiratory mortality at lag 1 is barely distinguishable from null
(\textbf{RR 1.010, 95\,\% CI 1.000--1.019}), but it separates cleanly on pollen load:
\textbf{1.021 (1.005--1.037)} above the 75th percentile against \textbf{1.003
(0.987--1.018)} below it, and the split is robust to the choice of cutoff. A mass-only
instrument in Fukuoka in March therefore returns a number whose health meaning changes
by a factor of seven depending on a biogenic co-pollutant it is blind to.
Çulpan et al.\ (\textit{Rev Environ Health}) reach the same conclusion by subtraction
across \textbf{21 studies and 990{,}157} children: the fine fraction is the only exposure
with a consistent blood-pressure association (\textbf{+0.59\,mmHg systolic per
10\,$\mu$g/m$^3$}, 0.28--0.90), the coarse-fraction estimates do not survive sensitivity
analysis, and NO$_2$ is null --- but the authors' own stated limitation is that they could
not fit multipollutant models because the primary studies did not report the exposure
covariance. Both papers are asking the measurement side for the same thing: a
co-exposure-resolved record, not a scalar.
\end{keybox}

\vspace{3mm}

% Bottom margin proofed: the exec brief closes clear of the folio at 2 baselines.
\enlargethispage{2\baselineskip}
% -------------------------------------------------------------- exec brief
\section*{\sffamily\bfseries\color{Deep}What changed today}
\vspace{-1.5mm}{\color{Amber}\rule{\linewidth}{1.1pt}}\vspace{2.5mm}

\footnotesize
\begin{itemize}
\item \textbf{The thinnest window this watch has run: 14 records against 26 yesterday
and 27 on 6 August.} This is a calendar artefact, not a signal --- 8 August is a Saturday
and PubMed's entry-date volume collapses at weekends. The PubMed connector returned
\textbf{9 hits on the health axis and zero on the instrumentation axis}, the third
consecutive empty sensing query from that leg, and the local harvester returned 4 and 0.
The corpus is held up almost entirely by \textbf{Crossref by ISSN, which supplied 6 of
the 7 sensing and exposure records}. Treat the daily counts for 8 August as a floor when
they enter the weekly trend line.

\item \textbf{Six of fourteen records --- and six of seven on the sensing axis --- carry no
abstract anywhere.} All six are \textit{Atmospheric Environment} / \textit{Atmospheric
Pollution Research} deposits made today; Crossref, Europe PMC and Semantic Scholar all
return metadata only, and Semantic Scholar has no record of the DOIs at all. They are
listed with design, metric and region but are \textbf{not summarised}, because a summary
written from a title is a fabrication. The highest-value pickups are Madalipay et al.\ on
provincial PM$_{2.5}$ conversion factors for South Korea and Cuerda Barcaiztegui et al.\
on a two-year classroom ventilation and IAQ field study in Madrid.

\item \textbf{Five effect estimates with intervals, down from sixteen yesterday}, and
three of them are the same Kyushu interaction reported at different pollen strata.
Two are co-pollutant (NO$_2$, indoor air nicotine). \textbf{No PM ratio estimate from a
cohort entered the corpus today.} Çulpan et al.'s meta-analytic estimates are mean
differences in mmHg and are deliberately excluded from the forest plot, which is a
log-ratio axis --- plotting 0.59 against a null of 1.0 would render a harmful effect as
protective. They are carried in the digest text and the register instead.

\item \textbf{The one high-temporal-resolution result is a co-pollutant one.} Shukla
et al.\ (\textit{BMJ Open}) fit a Poisson GAM/DLNM on \emph{hourly} NO$_2$ against
18{,}943 adult asthma admissions in Birmingham and Solihull, and locate the association
at a \textbf{3--6 hour lag} (RR 1.13, 1.01--1.26 at the 18\,$\mu$g/m$^3$ mean over
lag 0--24\,h, daytime). The lag structure is the transferable part: an exposure signal
resolvable at 3--6 hours is invisible to a 24-hour filter-based method and is exactly the
regime distributed low-cost sensing is built for. It is NO$_2$, not PM, so it enters as a
methodological rather than a substantive record.

\item \textbf{Nothing new on sensor calibration, drift or co-location.} Three Consensus
sweeps (calibration and drift, co-location against reference monitors, satellite--surface
fusion with humidity correction) returned 20 hits each, capped at 3 per sweep by the free
tier, and \textbf{every returned DOI resolved to 2019--2025 or to a record already in
\texttt{seen.json}}. OpenAlex returned HTTP 429 for the sixth consecutive run. The
calibration literature was not silent today; the available discovery paths simply cannot
see a single day of it.
\end{itemize}

\clearpage

% ==================================================================== ANALYTICS
\section*{\sffamily\bfseries\color{Deep}Corpus analytics}
\vspace{-1.5mm}{\color{Amber}\rule{\linewidth}{1.1pt}}\vspace{3mm}

\noindent\includegraphics[width=\linewidth]{\FIGDIR/f1_subtopics.png}
\figcap{Subtopic assignment of the 14 in-scope records, single-label by dominant
contribution. \textit{Burden, policy \& mitigation} and \textit{Exposure assessment \&
modelling} tie at 3; adding \textit{Sensing, forecasting \& instrumentation} (2) and
\textit{Occupational \& indoor} (2) puts \textbf{7 of 14 records on the measurement side}.
Three of the ten canonical clusters are empty today --- \textit{Neuro / mental health},
\textit{Reproductive \& developmental} and \textit{Other clinical endpoints} --- which at
$n=14$ is expected sampling variation, not a field-level absence.}

\vspace{2mm}
\noindent\includegraphics[width=\linewidth]{\FIGDIR/f2_design_endpoint.png}
\figcap{Left --- study architecture. \textbf{\textit{Metadata only} is the single largest
slice at 6 of 14 (43\%)}, the highest share this watch has recorded, and it is an honest
category rather than a residual: those records are listed but uncharacterised. No
observational cohort entered today. Right --- health endpoint. \textbf{Nine of fourteen
records carry \emph{no} human health endpoint} (six metadata-only sensing deposits, one
control-technology bench study, one curriculum evaluation, one bibliometric review), so
the endpoint panel describes five records and is not a field-level distribution.}

\vspace{2mm}
\noindent\includegraphics[width=\linewidth]{\FIGDIR/f3_metric_geography.png}
\figcap{Left --- particle metric. \textbf{Ten of fourteen records sit in
\textit{PM (unspeciated) / emissions}}, which is a direct consequence of the
metadata-only share: without an abstract there is no basis for assigning a size cut.
Only one record (Samburova et al., wildfire-influenced urban PM$_{2.5}$ in northern
Nevada) is unambiguously PM$_{2.5}$-only, and \textbf{no record today measures ultrafine
particles} --- a gap now standing at six consecutive issues. Right --- geographic
provenance. North America contributes 5; South Asia and Latin America are absent again.}

\vspace{2mm}
\noindent\includegraphics[width=\linewidth]{\FIGDIR/f6_lifecourse.png}
\figcap{Life-course windows addressed. Counts are non-exclusive and total 7, not 14,
because the six metadata-only records and the curriculum evaluation address no life
stage. The paediatric blood-pressure meta-analysis spans childhood and adolescence and is
counted in both.}

\vspace{2mm}
\noindent\includegraphics[width=\linewidth]{\FIGDIR/f4_heatmap.png}
\figcap{Subtopic $\times$ study-architecture cross-tabulation. The structural feature
worth naming is that the entire \textit{Metadata only} column (6 records) lands in the
\textit{Exposure assessment}, \textit{Sensing} and \textit{Occupational \& indoor} rows
--- \textbf{the abstract-availability failure is not randomly distributed across the
corpus, it falls almost entirely on the instrumentation axis}. This is the fourth issue in
a row where the day's most methodologically relevant content is the least characterisable,
and it is a property of Elsevier's atmospheric-sciences deposit pipeline, not of the
literature.}

\vspace{2mm}
\noindent\includegraphics[width=\linewidth]{\FIGDIR/f5_forest.png}
\figcap{The five interval estimates in today's corpus, on a common logarithmic axis.
\textbf{Three of the five are the same Kyushu SPM--mortality association reported at
different pollen strata} and are therefore not independent; they are plotted together
because the contrast between them is the finding. Two are co-pollutant rather than PM
(Shukla, hourly NO$_2$; Tillery, indoor air nicotine). Exposure contrasts are
heterogeneous --- a daily SPM increment, an 18\,$\mu$g/m$^3$ hourly NO$_2$ mean, and a
binary detection indicator --- so the plot shows precision and direction only.
\textbf{Two estimates are deliberately absent}: Çulpan et al.'s mmHg mean differences,
which have no place on a ratio axis, and Tillery et al.'s long-COVID OR 3.69
(1.11--16.6), whose interval spans a factor of 15 and would compress every other estimate
to a point (same exclusion rule applied to Jiang et al.\ on 7 August). Both are given in
the digest text.}

\clearpage

% ==================================================================== DIGEST
\section*{\sffamily\bfseries\color{Deep}Paper-level digest}
\vspace{-1.5mm}{\color{Amber}\rule{\linewidth}{1.1pt}}\vspace{1mm}

{\fontsize{7.4}{9}\selectfont\color{Slate}Ordered by cluster. Each entry gives the
methodological core and the finding that matters, not an abstract paraphrase. Records
marked \textsc{metadata only} were surfaced without an abstract in Crossref, Europe PMC or
Semantic Scholar and are listed, not summarised.\par}

\band{Respiratory \& allergic\hfill 2 records}{Sage}

\paperentry{Hur et al.\ 2026}
{Pollen modifies the SPM--mortality association across eight Kyushu cities}
{Environ Pollut}
{Two-stage time-series over 262{,}288 all-cause deaths, eight cities, pollen seasons
(Feb--Apr) 1989--2014, with an explicit SPM\,$\times$\,pollen interaction term in stage one
and random-effects pooling in stage two. Respiratory mortality at lag 1:
\textbf{RR 1.010 (1.000--1.019)} overall, \textbf{1.021 (1.005--1.037)} on days above the
75th-percentile pollen concentration against \textbf{1.003 (0.987--1.018)} below,
stable across alternative cutoffs. The threat to inference is exposure misclassification
in the modifier, not the exposure: pollen is measured at a single station per city and
the interaction is fitted on a binary split of a highly skewed distribution, so the
contrast is a lower bound on a continuous effect. Suspended particulate matter is a
Japanese regulatory metric with a 100\,\% cut near 10\,$\mu$m and is not
interchangeable with PM$_{2.5}$. For distributed sensing the implication is concrete:
an optical counter with a fluorescence or coarse-mode channel can resolve the modifier and
the exposure on the same node, which no regulatory network in this study did.}
{10.1016/j.envpol.2026.128933}

\paperentry{Shukla et al.\ 2026}
{Hourly NO$_2$ and emergency admissions for acute asthma, Birmingham and Solihull}
{BMJ Open}
{Poisson generalised additive model with a distributed-lag non-linear term on
\emph{hourly} counts --- 18{,}943 adult asthma admissions to four acute hospitals,
June 2016--May 2022, linked through the PIONEER data hub. The mean 18\,$\mu$g/m$^3$
NO$_2$ exposure over lag 0--24\,h carried \textbf{RR 1.13 (1.01--1.26)} for daytime
admissions, with the association concentrated at a \textbf{3--6 hour lag} and a steeper
gradient in the most deprived quintile. Exposure is assigned from fixed monitors to
residence, so the hourly resolution is nominal rather than personal --- time-activity
displacement over a 3--6\,h window is precisely the scale at which residential assignment
fails, which biases toward the null. The pollutant is NO$_2$, so this enters the watch on
method rather than on substance: it is the cleanest recent demonstration that the
health-relevant exposure signal lives below the 24-hour averaging period that
filter-based reference methods impose.}
{10.1136/bmjopen-2025-110612}

\band{Cardiovascular \& metabolic\hfill 1 record}{Teal}

\paperentry{Çulpan et al.\ 2026}
{Long-term air pollution and blood pressure in children: updated systematic review and meta-analysis}
{Rev Environ Health}
{PRISMA-guided update over prior reviews plus PubMed and Web of Science to 27 March 2025;
21 studies, 990{,}157 participants in the qualitative synthesis, with meta-analysis
restricted to the fine and coarse fractions and NO$_2$ because cohort data for SO$_2$,
O$_3$ and CO were too sparse. \textbf{Fine PM: +0.59\,mmHg systolic per 10\,$\mu$g/m$^3$
(0.28--0.90)}, the only estimate consistent across sensitivity analyses. The coarse
fraction gave nominally tighter systolic (+0.28, 0.25--0.31) and diastolic (+0.27,
0.24--0.30) associations that \emph{did not survive} sensitivity analysis --- an
implausibly narrow interval for a pooled estimate, which usually indicates a single
dominant study. NO$_2$ was null for both. Several pollutants raised the odds of
hypertension as a dichotomous outcome. Two limitations are load-bearing: exposure
increments and hypertension definitions were not harmonised across primary studies, and
no multipollutant model could be fitted because the primary studies did not report
exposure covariance, so the pollutant-specific attribution is weaker than the point
estimates suggest. Note also that the PubMed deposit strips subscripts from every
particle-size label, so fraction assignment here is inferred from the stated
meta-analytic scope.}
{10.1515/reveh-2026-0032}

\band{Occupational \& indoor exposure\hfill 2 records}{Clay}

\paperentry{Tillery et al.\ 2026}
{Indoor air nicotine, NO$_2$ and COVID-19 outcomes in a Great Plains tribal community}
{Indoor Air}
{COVID-19 Wayakta He? study: 290 multigenerational Cheyenne River Sioux Tribe households,
each pair of adults at least 9 years apart in age, with one week of passive air-nicotine
and NO$_2$ badges plus saliva serology. Geometric mean air nicotine
0.052\,$\mu$g/m$^3$ (detected in 40.1\,\% of homes); geometric mean NO$_2$ 10.5\,ppb,
range 0.65--138, consistent with propane cooking and heating. \textbf{OR 1.21
(1.01--1.51)} for prior infection with detectable air nicotine; \textbf{OR 3.69
(1.11--16.6)} for self-reported long COVID with NO$_2$ exposure. The second interval
spans a factor of fifteen and rests on a self-reported outcome in a cross-sectional
design with one week of exposure measurement standing in for a multi-year window ---
it is a direction, not a magnitude. The study's real contribution is that it measured
the microenvironment at all: no ambient network in the Northern Plains resolves
household combustion, and community-deployed passive samplers are the only instrument
that reached it.}
{10.1155/ina/7749725}

\paperentry{Cuerda Barcaiztegui et al.\ 2026}
{Classroom ventilation management and user activities vs.\ indoor air quality, Madrid primary schools}
{Atmos Pollut Res}
{\textsc{metadata only} --- no abstract in Crossref, Europe PMC or Semantic Scholar at
the time of this issue. A two-year classroom field study on ventilation management and
occupant activity is directly on this watch's deployment axis and is the single highest-priority abstract re-check for the next run.}
{10.1016/j.apr.2026.103170}

% Proof 1 spilled 5 lines of this entry onto an orphan p.6. Per the 2026-08-06/07 log,
% \enlargethispage does not recover a spill this large and prose trimming reflows without
% converging; moving a whole band across the break does. The Mechanistic band now opens
% p.6 instead of closing p.5.
\clearpage

\enlargethispage{4\baselineskip}
\band{Mechanistic toxicology\hfill 1 record}{Violet}

\paperentry{Yibcharoenporn et al.\ 2026}
{Metabolic profiling of PM-induced developmental toxicity in zebrafish embryos}
{ACS Omega}
{Zebrafish embryos exposed 4--50\,hpf to fine and coarse PM at 500 and
1000\,$\mu$g/mL, with untargeted LC-MS/MS metabolomics. The fine fraction alone produced
a \textbf{22\,h hatching delay} and disrupted pantothenate/CoA, retinol and riboflavin
metabolism; riboflavin fell significantly at the higher dose. Glutathione responses
diverged by size --- reduced glutathione depleted under the fine fraction but
\emph{elevated} under the coarse --- which is the most interesting result here because it
argues against a single shared oxidative mechanism scaled by mass. The dose is the
problem: 500--1000\,$\mu$g/mL of suspended PM is three to four orders of magnitude above
any inhalable ambient concentration, so this establishes a mechanism, not a risk. The
authors' closing suggestion of vitamin supplementation is not supported by anything in
the design.}
{10.1021/acsomega.6c00972}

\band{Sensing, forecasting \& instrumentation\hfill 2 records}{Amber}

\paperentry{Samburova et al.\ 2026}
{Effects of western US wildfire smoke on urban air quality in northern Nevada}
{Atmos Pollut Res}
{\textsc{metadata only} --- no abstract available in Crossref, Europe PMC or Semantic
Scholar. Long-range transported wildfire smoke over an urban receptor is the same
problem class as Jin et al.\ (\textit{ACP}, 7 August) on aged smoke in Missoula; worth a
re-check once the publisher deposits the abstract.}
{10.1016/j.apr.2026.103168}

\paperentry{Astaneh et al.\ 2026}
{Source-tagged PM tracers in random-forest bias correction of ozone design-value predictions}
{Atmos Environ}
{\textsc{metadata only}. Using source-tagged particulate tracers as features in an ML
bias-correction layer over a photochemical model is a transferable construction even
though the target here is ozone, not PM.}
{10.1016/j.atmosenv.2026.122288}

\band{Exposure assessment \& modelling\hfill 3 records}{Amber}

\paperentry{Madalipay et al.\ 2026}
{Establishing PM$_{2.5}$ conversion factors for South Korea at provincial level}
{Atmos Pollut Res}
{\textsc{metadata only}. Provincial PM$_{2.5}$ conversion factors bear directly on how a
sparse regulatory network is extrapolated to unmonitored areas, which is the same
inference a low-cost network is deployed to replace. Priority re-check alongside the
Madrid classroom study.}
{10.1016/j.apr.2026.103171}

\paperentry{McCerery et al.\ 2026}
{Arctic transport of North American and Canadian wildfire aerosols}
{Atmos Environ}
{\textsc{metadata only}. The only record in today's corpus with a polar or remote-marine
receptor; long-range transport to a low-background site is where instrument detection
limits, not calibration slope, become the binding constraint.}
{10.1016/j.atmosenv.2026.122289}

\paperentry{Raptaki et al.\ 2026}
{Atmospheric BTEX in Heraklion, Crete: levels, sources and contribution to secondary aerosol formation}
{Atmos Environ}
{\textsc{metadata only}. Enters on the secondary-organic-aerosol precursor link rather
than on any direct PM measurement; tier C.}
{10.1016/j.atmosenv.2026.122283}

\band{Burden, policy \& mitigation\hfill 3 records}{Coral}

\paperentry{Gu et al.\ 2026}
{Electro-ultrasonic coupled agglomeration for removal of fine particles from biomass combustion}
{J Hazard Mater}
{Segmented continuous-flow bench platform combining bipolar charging ($\pm$16\,kV), a
300\,V/400\,Hz AC field and 16\,kHz/142\,dB ultrasound, with particle-trajectory
simulation and SEM. The coupled field reached \textbf{86.4\,\% light transmittance},
against single-field results 33.1 and 67.2 percentage points lower, with the extinction
coefficient down to 0.9\,m$^{-1}$; simulated cumulative particle displacement rose
$>$445\,\% as linear trajectories became three-dimensional helical paths, and SEM
confirmed 15--20\,$\mu$m agglomerates. The performance metric is the weakness: light
transmittance and extinction coefficient are \emph{optical} proxies, not gravimetric or
number-resolved removal efficiency, so the headline ``removal efficiency'' is not
directly comparable with any filtration literature that reports mass penetration. The
submicron fraction claimed as the target is precisely the fraction an extinction
measurement is least sensitive to.}
{10.1016/j.jhazmat.2026.143130}

\paperentry{Wojcik et al.\ 2026}
{An air pollution curriculum for pulmonary and critical care fellows}
{ATS Sch}
{Kern six-step curriculum development across five domains --- composition, health effects,
\emph{monitoring}, risk mitigation and environmental justice --- with a 20-item knowledge
scale built by modified Delphi and delivered as pre/post assessment over one academic
year with spaced repetition. Knowledge rose in all five domains ($p<0.001$). Single
programme, no control arm, self-selected respondents and a knowledge rather than
behaviour endpoint, so the effect size is uninterpretable; it is included because the
explicit inclusion of \emph{monitoring} as a clinical competency domain is a small but
real indicator of where sensor-derived exposure data is expected to land clinically.}
{10.1093/atsscholar/aapag039}

\paperentry{Kokol et al.\ 2026}
{Artificial intelligence in the management of climate-change-related health issues in Africa}
{Res Sq (preprint)}
{Scoping review over Scopus to October 2025, VOSviewer clustering a global against an
African literature corpus. The divergence is the finding: global work clusters on
environmental modelling, air-pollution mortality tracking and exposome integration,
African work on supervised classification for infectious-disease triage and food
security, with a ``compute deficit'', data silos and power instability named as barriers.
This is bibliometrics, not evidence --- a VOSviewer map measures what Scopus indexed, and
Scopus under-indexes African journals, so the gap is partly an artefact of the
instrument. Preprint, not peer reviewed.}
{10.21203/rs.3.rs-9910414/v1}

\clearpage

% ==================================================================== TABLE
\section*{\sffamily\bfseries\color{Deep}Corpus register}
\vspace{-1.5mm}{\color{Amber}\rule{\linewidth}{1.1pt}}\vspace{2.5mm}

{\fontsize{6.7}{8.4}\selectfont
\setlength{\tabcolsep}{2.4pt}
\begin{longtable}{@{}>{\RaggedRight}p{25mm}>{\RaggedRight}p{41mm}>{\RaggedRight}p{27mm}%
>{\RaggedRight}p{23mm}>{\RaggedRight}p{18mm}>{\RaggedRight}p{15mm}@{}}
\toprule
\textsf{\textbf{First author}} & \textsf{\textbf{Focus}} & \textsf{\textbf{Journal}} &
\textsf{\textbf{Design}} & \textsf{\textbf{Metric}} & \textsf{\textbf{Region}}\\
\midrule
\endfirsthead
\toprule
\textsf{\textbf{First author}} & \textsf{\textbf{Focus}} & \textsf{\textbf{Journal}} &
\textsf{\textbf{Design}} & \textsf{\textbf{Metric}} & \textsf{\textbf{Region}}\\
\midrule
\endhead
\multicolumn{6}{@{}l@{}}{\cellcolor{Mist}\textsf{\textbf{\textcolor{Deep}{Neuro / mental health}}}}\\[0.6mm]
\href{https://doi.org/10.1001/jamanetworkopen.2026.25015}{Dales} & Wildfire-sourced PM$_{2.5}$ and ED visits for migraine and other primary headache & \textit{JAMA Netw Open} & Case-crossover & PM$_{2.5}$ (wildfire) & Canada\\
\href{https://doi.org/10.1093/ageing/afag215}{Glover} & Umbrella meta-review of urban environmental exposures and cognitive health & \textit{Age Ageing} & Umbrella review / meta-analysis & PM$_{2.5}$ & Global\\
\href{https://doi.org/10.1016/j.nsa.2026.107019}{Goossens} & In utero air pollution exposure and pre-/postnatal brain development & \textit{Neurosci Appl} & Narrative review & PM$_{2.5}$ / PM$_{10}$ / PM$_{\mathrm{coarse}}$ & Global\\
\href{https://doi.org/10.1016/j.isci.2026.116617}{Ma} & PM$_{2.5}$-induced neurotoxicity and depression: exposure risks and mechanisms & \textit{iScience} & Narrative review & PM$_{2.5}$ + components & Global\\
\href{https://doi.org/10.3760/cma.j.cn112150-20250804-00752}{Pan} & Maternal PM exposure and offspring neurodevelopment: effects and mechanisms & \textit{Zhonghua Yu Fang Yi Xue Za Zhi} & Narrative review & PM (bulk) & China\\
\addlinespace[0.9mm]
\multicolumn{6}{@{}l@{}}{\cellcolor{Mist}\textsf{\textbf{\textcolor{Teal}{Cardiovascular \& metabolic}}}}\\[0.6mm]
\href{https://doi.org/10.3760/cma.j.cn112338-20251120-00843}{Huang} & PM$_{2.5}$ and dyslipidaemia in rural Chinese children; dietary effect modification & \textit{Zhonghua Liu Xing Bing Xue Za Zhi} & Cross-sectional (multistage) & PM$_{2.5}$ & China\\
\href{https://doi.org/10.32394/pe/218966}{Kazmierska} & Air pollution and cardiovascular disease: review of current data & \textit{Przegl Epidemiol} & Narrative review & PM$_{2.5}$ / PM$_{10}$ & Global\\
\href{https://doi.org/10.1007/s00125-026-06814-2}{Xi} & Air pollution, built environment and temperature vs. prediabetes phenotypes & \textit{Diabetologia} & Prospective cohort & PM$_{2.5}$ & Germany\\
\href{https://doi.org/10.1161/JAHA.125.046917}{Yang} & Short-term PM$_{2.5}$ and hospitalisation/ED use in atrial fibrillation & \textit{J Am Heart Assoc} & Case-crossover & PM$_{2.5}$ (wildfire-influenced) & USA\\
\addlinespace[0.9mm]
\multicolumn{6}{@{}l@{}}{\cellcolor{Mist}\textsf{\textbf{\textcolor{Sky}{Reproductive \& developmental}}}}\\[0.6mm]
\href{https://doi.org/10.1016/j.jes.2025.10.043}{Chen} & Prenatal air pollution, fetal growth trajectory phenotypes and child neurodevelopment & \textit{J Environ Sci (China)} & Prospective cohort & PM$_{2.5}$ / PM$_{10}$ & China\\
\href{https://doi.org/10.1016/j.ecoenv.2026.120566}{Huang} & Joint prenatal bisphenol + air-pollutant mixture exposure and preschool IQ & \textit{Ecotoxicol Environ Saf} & Prospective cohort (mixtures) & PM$_{2.5}$ / PM$_{10}$ & China\\
\href{https://doi.org/10.1007/s40572-026-00552-8}{Peterson} & Ambient air pollution, PAHs and pubertal development: critical review & \textit{Curr Environ Health Rep} & Critical review & PM$_{2.5}$ components & Global\\
\href{https://doi.org/10.1016/j.envres.2026.125313}{Xie} & Long-term PM$_{2.5}$ and timing of pubertal development in boys and girls & \textit{Environ Res} & Prospective cohort & PM$_{2.5}$ & China\\
\addlinespace[0.9mm]
\multicolumn{6}{@{}l@{}}{\cellcolor{Mist}\textsf{\textbf{\textcolor{Sage}{Respiratory \& allergic}}}}\\[0.6mm]
\href{https://doi.org/10.1007/s00484-026-03289-z}{Bayar Muluk} & Climate change and otolaryngology: airway disease impacts and emission reduction & \textit{Int J Biometeorol} & Narrative review & PM$_{2.5}$ & Global\\
\href{https://doi.org/10.4103/mgr.MEDGASRES-D-25-00164}{Liao} & Short-term air pollution, respiratory symptoms and lung function in asthma & \textit{Med Gas Res} & Panel study & PM$_{2.5}$ / PM$_{10}$ & China\\
\href{https://doi.org/10.3390/epidemiologia7040097}{Sanchez de Prada} & Ambient PM and asthma-coded admissions in children 0-3 y, Spain & \textit{Epidemiologia} & Ecological panel & PM$_{2.5}$ / PM$_{10}$ & Spain\\
\href{https://doi.org/10.1016/j.jhazmat.2026.143076}{Yang} & Allergenic vs. toxic risk of airborne seed hairs compared with ambient PM$_{2.5}$ & \textit{J Hazard Mater} & Field measurement & PM$_{2.5}$ + bioaerosol & China\\
\addlinespace[0.9mm]
\multicolumn{6}{@{}l@{}}{\cellcolor{Mist}\textsf{\textbf{\textcolor{Violet}{Mechanistic toxicology}}}}\\[0.6mm]
\href{https://doi.org/10.1016/j.jes.2025.12.074}{Ding} & PM$_{2.5}$ promotes vascular calcification via miR-27b-5p/HES1 signalling & \textit{J Environ Sci (China)} & In vitro + organoid + murine & PM$_{2.5}$ & China\\
\href{https://doi.org/10.1007/s00360-026-01698-5}{Maraschi} & Osmoregulatory disruption and metal bioaccumulation from settleable atmospheric PM & \textit{J Comp Physiol B} & Controlled exposure (ecotox) & Settleable APM & Brazil\\
\href{https://doi.org/10.1016/j.jes.2025.11.018}{Wu} & Macrophage NLRP3 inflammasome in real-ambient PM-induced hepatic glucose dysregulation & \textit{J Environ Sci (China)} & Real-ambient murine + co-culture & Real-ambient PM & China\\
\addlinespace[0.9mm]
\multicolumn{6}{@{}l@{}}{\cellcolor{Mist}\textsf{\textbf{\textcolor{Clay}{Occupational \& indoor}}}}\\[0.6mm]
\href{https://doi.org/10.3389/fmedt.2026.1796828}{Chen} & Nested suction instrument to cut intraoperative surgical-smoke PM$_{2.5}$ & \textit{Front Med Technol} & Bench / device experiment & PM$_{2.5}$ (surgical smoke) & China\\
\href{https://doi.org/10.1093/annweh/wxag055}{Cole} & Determinants of respirable crystalline silica exposure in Australian tunnelling & \textit{Ann Work Expo Health} & Observational exposure study & Respirable dust / RCS & Australia\\
\href{https://doi.org/10.1016/j.jhazmat.2026.143081}{Ding} & Stage-resolved hair Zn-isotope dynamics as a firefighting combustion-exposure tracer & \textit{J Hazard Mater} & Repeated-measures biomarker & Combustion-derived PM & China\\
\href{https://doi.org/10.3390/jox16040119}{Ramsden} & Xenobiotic hazards in aircraft cabin air & \textit{J Xenobiot} & Review + risk economics & Ultrafine particles & Global\\
\addlinespace[0.9mm]
\multicolumn{6}{@{}l@{}}{\cellcolor{Mist}\textsf{\textbf{\textcolor{Amber}{Sensing, forecasting \& instrumentation}}}}\\[0.6mm]
\href{https://doi.org/10.1016/j.jenvman.2026.130515}{Shi} & Adaptive spatiotemporal graph transformer for multi-site multi-step PM$_{2.5}$ forecasting & \textit{J Environ Manage} & Deep-learning model & PM$_{2.5}$ & China / India\\
\href{https://doi.org/10.3390/bios16070380}{Tsuchiya} & Seasonal and indoor-outdoor characteristics of fluorescent aerosol particles in offices & \textit{Biosensors} & Real-time sensor campaign & Bioaerosol / coarse & Japan\\
\addlinespace[0.9mm]
\multicolumn{6}{@{}l@{}}{\cellcolor{Mist}\textsf{\textbf{\textcolor{Amber}{Exposure assessment \& modelling}}}}\\[0.6mm]
\href{https://doi.org/10.1029/2025GH001621}{Solorzano-Araque} & Multiscale exposure modelling: LOTOS-EUROS + mobility surveys + particle toxicity & \textit{GeoHealth} & CTM + mobility integration & PM$_{2.5}$ / PM$_{10}$ & Colombia\\
\addlinespace[0.9mm]
\multicolumn{6}{@{}l@{}}{\cellcolor{Mist}\textsf{\textbf{\textcolor{Coral}{Burden, policy \& mitigation}}}}\\[0.6mm]
\href{https://doi.org/10.1002/hsr2.72886}{Karamzad} & Ischaemic heart disease burden in MENA 1990-2021; PM attributable fraction 28.2\% & \textit{Health Sci Rep} & GBD secondary analysis & Ambient PM & MENA\\
\href{https://doi.org/10.1371/journal.pone.0347989}{Rees} & HEPA filtration to reduce respiratory infection in care homes & \textit{PLoS One} & RCT process evaluation & Filtered PM & UK\\
\href{https://doi.org/10.1016/j.jes.2026.01.009}{Shan} & Policy-driven emission reductions and post-pandemic rebound, high-resolution inventories & \textit{J Environ Sci (China)} & Emission inventory & PM emissions & China\\
\addlinespace[0.9mm]
\multicolumn{6}{@{}l@{}}{\cellcolor{Mist}\textsf{\textbf{\textcolor{Slate}{Other clinical endpoints}}}}\\[0.6mm]
\href{https://doi.org/10.3390/dj14070440}{Lupita} & PM$_{2.5}$, preventive-care capacity and caries/edentulism burden across the EU & \textit{Dent J} & Ecological panel & PM$_{2.5}$ & EU-27\\
\href{https://doi.org/10.1016/j.jdin.2026.06.008}{Park} & Long-term PM exposure and risk of androgenetic alopecia & \textit{JAAD Int} & Nationwide cohort & PM$_{2.5}$ / PM$_{10}$ & South Korea\\
\href{https://doi.org/10.3171/2026.3.PEDS2615}{Sun} & Environmental exposure and area deprivation vs. paediatric hydrocephalus outcomes & \textit{J Neurosurg Pediatr} & Registry cohort + geospatial & PM$_{2.5}$ (satellite-derived) & USA\\
\addlinespace[0.9mm]
\bottomrule
\end{longtable}}


\vspace{4mm}
\begin{keybox}
{\sffamily\bfseries\footnotesize\color{Deep}Provenance and method}\\[1.2mm]
{\fontsize{7.2}{8.8}\selectfont
Entry window \textbf{2026-08-08 to 2026-08-08}, contiguous with the 7 August issue; no
gap. Legs run: \textbf{PubMed} connector, two independent queries against
\texttt{[Date - Entry]} --- health axis 9 hits, instrumentation axis \textbf{0}; a third
MeSH-only sweep returned 10, of which 9 were water, soil or pharmaceutical-cleanroom
monitoring and were rejected. The local harvester returned 4 health and 0 sensing.
\textbf{Europe PMC} on \texttt{CREATION\_DATE} returned 2 preprints, 1 carried.
\textbf{Crossref by ISSN} returned 7 and is the only productive sensing leg, 6 carried.
\textbf{Consensus} $\times$3 sweeps returned 20 each, capped at 3 per sweep by the free
tier; all 9 returned records resolved to 2019--2025 DOIs or to \texttt{seen.json}, 0
carried. \textbf{arXiv} 50 records, none published or updated on or after 1 August.
\textbf{OpenAlex} HTTP 429 (\texttt{from\_created\_date} leg remains broken).
\textbf{13 records rejected}, each with a reason in \texttt{state/rejected.jsonl}.
Six carried records have no abstract in any of the three checked sources and are
listed without a summary; they are marked \textsc{metadata only} in the digest and in
the architecture donut.
Subtopic, design, particle-metric and geography labels are assigned from title, abstract,
keywords and MeSH; assignment is single-label by dominant contribution and is a
judgement, not a controlled vocabulary.
All DOI links resolve to the publisher of record. Content derived from PubMed and
Crossref metadata; abstracts remain the copyright of their respective publishers.\par}
\end{keybox}

\end{document}
